\sectionkicker{Source-backed technical notes}
\subsection*{生成から行動へ――Plugin、Function Calling、Tool UseがAPI層を作った: Technical Notes}
本欄は記事本文で圧縮した一次資料上の情報を、比較・再検証しやすい形で整理したものである。時系列、確認済みの主張、留意点、一次資料URLを掲載する。
\smallskip
\noindent{\small\textit{Technical Notesでは、一次資料から確認した公開・提供・研究上の事実を記載する。数値・能力評価や提供元・著者による比較表現は、独立再現済みの結果ではなく帰属claimとして扱う。}}\par
\medskip
\subsection*{Theme at a glance}
\addcontentsline{toc}{subsection}{Theme at a glance}
\begin{center}
\footnotesize
\begin{tabularx}{\linewidth}{@{}p{0.20\linewidth}p{0.10\linewidth}p{0.14\linewidth}X@{}}
\toprule
資料 & 位置づけ & 種別 & 時系列 \\
\midrule
OpenAI tool/action surface 2023 & 主要資料 & API & 2023-03-23, 2023-06-13, 2023-11-06 \\
Toolformer & 補足資料 & 論文 & 2023-02-09 \\
\bottomrule
\end{tabularx}
\end{center}
\normalsize

\begin{technicalnote}{OpenAI tool/action surface 2023}{主要資料}
\begingroup
% reader-facing Technical Notes break policy
\widowpenalty=10000
\clubpenalty=10000
\displaywidowpenalty=10000
\begin{tabularx}{\linewidth}{@{}>{\bfseries}p{0.22\linewidth}X@{}}
組織 & OpenAI \\
種別 & API \\
時系列 & 2023-03-23, 2023-06-13, 2023-11-06 \\
\end{tabularx}
\smallskip
{\bfseries 一次資料から整理したtechnical points}
\begin{itemize}[leftmargin=1.5em,itemsep=0.35em]
\item \textbf{一次情報で確認できる事実}: OpenAIの選定済み一次資料では、2023年3月の ChatGPT plugins は ChatGPT から外部の最新情報・計算・第三者サービスへ接続する tool surface として initial support が公表された。これは後の API-level function calling とは別の ChatGPT product surface として扱う。 2023年6月の API update では function calling capability が API model に追加され、model output と developer-defined function / external tool の間を構造化する interface が提示された。ChatGPT plugin の user-facing rollout と同一の availability event とはみなさない。 2023年11月の GPTs は instructions・extra knowledge・skills を組み合わせた custom versions of ChatGPT を end user が構成する product surface として公表された。 同日の DevDay developer release では GPT-4 Turbo の 128K context、Assistants API、GPT-4 Turbo with Vision、DALL·E 3 API などが別々の developer product / model availability として発表されたため、GPT builder の product lifecycle と hosted API lifecycle を区別する。
\end{itemize}
\begin{minipage}{\linewidth}
% reader-facing Technical Notes generic boundary/source tail group
\begin{samepage}
{\bfseries 一次資料}
\begingroup\sloppy
\Urlmuskip=0mu plus 2mu\relax
\begin{itemize}[leftmargin=1.5em,itemsep=0.25em]
\item {\scriptsize\url{https://openai.com/index/chatgpt-plugins}}
\item {\scriptsize\url{https://openai.com/index/function-calling-and-other-api-updates}}
\item {\scriptsize\url{https://openai.com/index/introducing-gpts}}
\item {\scriptsize\url{https://openai.com/index/new-models-and-developer-products-announced-at-devday}}
\end{itemize}
\endgroup
\end{samepage}
\end{minipage}
\endgroup
\end{technicalnote}
\medskip

\begin{technicalnote}{Toolformer}{補足資料}
\begingroup
% reader-facing Technical Notes break policy
\widowpenalty=10000
\clubpenalty=10000
\displaywidowpenalty=10000
\begin{tabularx}{\linewidth}{@{}>{\bfseries}p{0.22\linewidth}X@{}}
組織 & authors \\
種別 & 論文 \\
時系列 & 2023-02-09 \\
\end{tabularx}
\smallskip
{\bfseries 一次資料から整理したtechnical points}
\begin{itemize}[leftmargin=1.5em,itemsep=0.35em]
\item \textbf{一次情報で確認できる事実}: authorsの選定済み一次資料では、Toolformer は language model 自身が external tools を simple APIs 経由で利用する方法を self-supervised に学習する枠組みとして提案され、少数の API-use demonstrations から学習データを拡張する。 著者らの方法では model がどの API を呼ぶか、いつ呼ぶか、どの引数を渡すか、返された結果を future token prediction にどう組み込むかを学習対象にする。 論文では calculator、question answering、search、translation、calendar など複数の tool classes を扱う。性能比較は原著者の実験結果として扱い、独立再現済みとはみなさない。
\end{itemize}
\begin{minipage}{\linewidth}
% reader-facing Technical Notes generic boundary/source tail group
\begin{samepage}
{\bfseries 一次資料}
\begingroup\sloppy
\Urlmuskip=0mu plus 2mu\relax
\begin{itemize}[leftmargin=1.5em,itemsep=0.25em]
\item {\scriptsize\url{http://arxiv.org/abs/2302.04761v1}}
\end{itemize}
\endgroup
\end{samepage}
\end{minipage}
\endgroup
\end{technicalnote}
\medskip
